\documentclass[11pt,a4paper]{article}

%% ---- Packages -------------------------------------------------------
\usepackage[utf8]{inputenc}
\usepackage[T1]{fontenc}
\usepackage{lmodern}
\usepackage[margin=2.5cm]{geometry}
\usepackage{amsmath,amssymb}
\usepackage{booktabs}
\usepackage{longtable}
\usepackage{array}
\usepackage{multirow}
\usepackage{graphicx}
\usepackage{xcolor}
\usepackage{hyperref}
\usepackage{url}
\usepackage{microtype}
\usepackage{setspace}
\usepackage{titlesec}
\usepackage{caption}
\usepackage{subcaption}
\usepackage{enumitem}
\usepackage{parskip}
\usepackage{natbib}
\usepackage{float}
\usepackage{tabularx}

%% ---- Hyperref setup -------------------------------------------------
\hypersetup{
  colorlinks = true,
  linkcolor  = blue!70!black,
  citecolor  = blue!70!black,
  urlcolor   = blue!70!black,
  pdftitle   = {Assessing the Effectiveness of Large Language Models as Polling Participants in Qualitative Research: A 2023-2026 Longitudinal Extension},
  pdfauthor  = {},
}

%% ---- Title ---------------------------------------------------------
\title{%
  \Large\bfseries
  Assessing the Effectiveness of Large Language Models as\\
  Polling Participants in Qualitative Research:\\
  A 2023--2026 Longitudinal Extension%
}
\author{%
  Tarun Gogineni\\
  \textit{Palo Alto Networks}\\
  \texttt{tarun.gogineni@paloaltonetworks.com}
}
\date{May 2026}

%% ====================================================================
\begin{document}
\maketitle
\thispagestyle{empty}

%% ---- Abstract -------------------------------------------------------
\begin{abstract}
Large language models (LLMs) have emerged as a promising tool for simulating
survey respondents in qualitative research, enabling low-cost, scalable pilot
studies without requiring human participant recruitment. This paper presents a
longitudinal extension of a 2023 study that evaluated GPT-3.5-Turbo and
LLaMA-2-7B as synthetic survey respondents. We introduce three new models ---
GPT-4o, Claude Sonnet 4-6, and LLaMA-3.3-70B --- and evaluate them across the
same 12-question software-developer survey ($n{=}314$ human baseline) under five
prompt variants, including a novel Chain-of-Thought (CoT) condition. We
additionally assess each model's scientific question-answering capability on 50
papers from the Qasper benchmark. Statistical analyses include Fleiss' Kappa
(inter-rater agreement), chi-square tests (distribution similarity), $t$-tests
(demographic sensitivity), Cram\'{e}r's~$V$ (effect size), and one-way ANOVA.

Our findings reveal a \textit{Kappa Paradox}: RLHF-aligned 2026 models exhibit
nearly double the within-group Fleiss' Kappa of their 2023 predecessors
(GPT-4o: $\kappa{=}0.480$; Claude Sonnet 4-6: $\kappa{=}0.487$;
LLaMA-3.3-70B: $\kappa{=}0.480$ versus GPT-3.5: $\kappa{=}0.241$, human:
$\kappa{=}0.103$), indicating dramatically increased response uniformity rather
than improved demographic diversity. Concurrently, chi-square tests show that
2026 model response distributions are substantially less divergent from the
reference distribution than their predecessors (GPT-4o: $\chi^2{=}2.86$,
$p{=}0.581$; Claude: $\chi^2{=}1.88$, $p{=}0.759$ versus 2023 baseline:
$\chi^2{=}18.75$, $p{<}0.001$), while all models --- including 2023 baselines
--- show zero statistically significant demographic sensitivity across Age,
Gender, and Experience dimensions. On the Qasper scientific QA benchmark, 2026
models substantially outperform 2023 baselines on BLEU-4 (LLaMA-3.3: 0.249
vs.\ LLaMA-2: 0.107), ROUGE-L (LLaMA-3.3: 0.395 vs.\ LLaMA-2: 0.282), and
METEOR (Claude: 0.523 vs.\ GPT-3.5: 0.298). These results have direct
implications for the use of LLMs as synthetic respondents in software
engineering qualitative research.

\bigskip
\noindent\textbf{Keywords:} large language models, survey simulation, synthetic
respondents, RLHF alignment, Fleiss' Kappa, Qasper, longitudinal study,
software engineering
\end{abstract}

\newpage
\tableofcontents
\newpage

%% ====================================================================
\section{Introduction}
\label{sec:intro}

Qualitative research in software engineering frequently relies on surveys and
interviews with human participants. However, recruiting diverse, representative
participants is costly, time-consuming, and often limited by self-selection bias.
The emergence of instruction-tuned LLMs has raised a compelling question: can
these models serve as \emph{synthetic respondents}, approximating human survey
behaviour well enough to pilot study instruments, validate questionnaire designs,
or augment small-sample studies?

A foundational 2023 study~\citep{original2023} explored this question by
prompting GPT-3.5-Turbo and LLaMA-2-7B with demographic profiles and measuring
whether their responses aligned with those of 314 human software developers
across a 12-question survey. The study found that both models exhibited higher
inter-rater agreement (Fleiss' Kappa $= 0.241$ for GPT-3.5, $0.238$ for
LLaMA-2) than human respondents ($\kappa = 0.103$), suggesting that LLMs
produce more internally consistent but less humanistically diverse responses.
Chi-square tests confirmed statistically significant divergence between model
and human response distributions ($\chi^2 = 18.75$, $p < 0.001$).

Since 2023, the LLM landscape has transformed dramatically. OpenAI's GPT-4o,
Anthropic's Claude Sonnet 4-6, and Meta's LLaMA-3.3-70B represent a new
generation of models trained with Reinforcement Learning from Human Feedback
(RLHF) at unprecedented scale. RLHF alignment has been shown to improve
instruction-following, reduce harmful outputs, and increase response
coherence~\citep{ouyang2022}. However, its effect on \emph{demographic
sensitivity} --- the ability of an LLM to produce meaningfully different
responses when prompted with different demographic profiles --- remains
underexplored. Closely related work on LLMs as survey participants in
cybersecurity and enterprise security contexts has also highlighted systematic
persona-alignment failures in RLHF-trained models~\citep{ssrn6829506}.

This paper addresses three research questions:

\begin{description}[leftmargin=2em]
  \item[\textbf{RQ1}] Do 2026 LLMs simulate human survey respondents more
    accurately than their 2023 predecessors, as measured by Fleiss' Kappa
    proximity to the human baseline?
  \item[\textbf{RQ2}] Do 2026 models produce higher-quality answers on the
    Qasper scientific QA benchmark, as measured by automatic and stylistic
    metrics?
  \item[\textbf{RQ3}] Did RLHF alignment reduce demographic sensitivity in LLM
    survey responses, as measured by $t$-test significance across Age, Gender,
    and Experience demographics?
\end{description}

\medskip\noindent
We make the following contributions:

\begin{enumerate}
  \item The first longitudinal comparison of LLM survey simulation capability
    across two model generations (2023 vs.\ 2026), using identical evaluation
    protocols.
  \item A new Chain-of-Thought prompt variant (P5) added to the existing four
    prompt templates (P1--P4), investigating whether structured reasoning
    improves demographic persona adherence.
  \item A comprehensive evaluation of three 2026 models on the Qasper benchmark
    using BLEU-1/2/3/4, ROUGE-1/2/L, METEOR, and five stylistic metrics
    (relevance, fluency, formality, readability, correctness).
  \item Open-source release of the complete inference pipeline, all model
    responses, and analysis scripts.
\end{enumerate}

%% ====================================================================
\section{Related Work}
\label{sec:related}

\subsection{LLMs as Survey Respondents}

The use of LLMs to simulate human survey responses has attracted growing
attention. \citet{argyle2023} introduced the concept of ``silicon sampling,''
demonstrating that GPT-3 could reproduce attitudinal patterns from the American
National Election Studies when prompted with demographic information. Their work
showed that LLM responses correlate with subgroup survey responses at rates
exceeding random baseline, though they cautioned against treating LLM-simulated
data as equivalent to human data.

\citet{santurkar2023} extended this analysis to GPT-3.5, finding that default
(non-demographically prompted) LLMs exhibit systematic opinion biases that skew
toward liberal, Western, and educated viewpoints. Prompted models showed
improved demographic sensitivity, though significant gaps remained for
underrepresented groups.

In the software engineering domain, the 2023 study that this paper
extends~\citep{original2023} was among the first to apply LLM survey simulation
specifically to developer-focused questionnaires, validating results against a
real $n{=}314$ developer survey. The SSRN working paper by \citet{ssrn6829506}
further examines LLM persona fidelity in professional and enterprise contexts,
providing complementary evidence for the profile-blindness phenomenon we
document here.

\subsection{RLHF Alignment and Response Diversity}

Reinforcement Learning from Human Feedback~\citep{ouyang2022} has become the
dominant fine-tuning paradigm for instruction-following LLMs. By training models
to produce responses rated highly by human raters, RLHF improves helpfulness
and safety but also introduces a tendency toward ``safe,'' high-consensus
answers. \citet{perez2022} demonstrated that RLHF-aligned models show reduced
response variance across equivalent prompts, which has implications for their
use as synthetic respondents: a model that always defaults to the most socially
acceptable answer may appear unbiased but actually underrepresents minority
viewpoints.

\subsection{Prompt Engineering for Persona Simulation}

The influence of prompt formulation on LLM demographic simulation has been
systematically explored by several groups. \citet{argyle2023} found that
explicit demographic backstory prompts substantially improved persona alignment.
The 2023 study~\citep{original2023} compared four prompt styles --- expert NLP
persona~(P1), plain QA~(P2), concise instruction~(P3), and novice/child
persona~(P4) --- and found P4 (novice/child) to yield the highest alignment with
human distributions. We extend this with a Chain-of-Thought (P5) prompt,
hypothesising that step-by-step reasoning may reduce pattern-matching shortcuts.

\subsection{Scientific Question Answering Benchmarks}

The Qasper dataset~\citep{dasigi2021} provides 1,585 questions over 888 NLP
research papers, with extractive and free-form gold answers from human
annotators. It has been widely used to evaluate document-level reading
comprehension. We use a 50-paper random sample (seed\,=\,42) from the test
split, consistent with prior work, to evaluate factual QA capability across
model generations.

\subsection{Evaluation Metrics}

BLEU~\citep{papineni2002}, ROUGE~\citep{lin2004}, and METEOR~\citep{banerjee2005}
remain standard automatic metrics for text generation evaluation. For stylistic
evaluation, we employ TF-IDF cosine similarity (relevance), type-token ratio
(lexical diversity / fluency proxy), MTLD (formality), Flesch Reading Ease
(readability), and FuzzyWuzzy token-set ratio (correctness). Due to system
memory constraints, we use TF-IDF cosine similarity as a proxy for
BERTScore~\citep{zhang2019} in the 2026 evaluation; 2023 BERTScore values from
the original paper used the \texttt{microsoft/deberta-xlarge-mnli} model and
are not directly comparable.

%% ====================================================================
\section{Methodology}
\label{sec:method}

\subsection{Human Survey Baseline}

The human baseline consists of $n{=}314$ software developer responses to a
12-question survey administered at a university computing department.
Respondents ranged from undergraduate students to professional developers, with
demographic distributions across Age (18--22, 23--26, 27--35 years), Gender
(Man, Woman), and Experience ($<$\,1 year, 1--3 years, 3--5 years). Questions
covered development environment preferences, learning strategies, professional
challenges, programming language selection, collaboration, innovation versus
deadlines, societal contributions, and ethical decision-making scenarios.

\subsection{Questionnaire Stimulus}

The 12-question survey instrument from the original 2023 study is reproduced
verbatim. Questions include:

\begin{enumerate}[noitemsep]
  \item Preferred development environment (Windows / macOS / Linux / Other)
  \item How do you learn to code? (multi-select)
  \item Biggest challenge as a developer (4 options)
  \item Programming language selection criteria (4 options)
  \item Team communication and collaboration (4 options)
  \item Staying up-to-date with industry changes (multi-select)
  \item Balancing innovation vs.\ deadlines (4 options)
  \item Software's contribution to societal challenges (multi-select)
  \item AI-driven employee dismissal ethical scenario (3 options)
  \item Unfamiliar language under deadline scenario (4 options)
  \item Critical pre-release bug scenario (4 options)
  \item SaaS delivery with critical bug scenario (5 options)
\end{enumerate}

\subsection{Demographic Profiles}

Ten synthetic profiles are constructed by crossing Age $\times$ Gender
$\times$ Ethnicity $\times$ Education $\times$ Experience
(Table~\ref{tab:profiles}).

\begin{table}[H]
\centering
\caption{Synthetic Demographic Profiles}
\label{tab:profiles}
\small
\begin{tabular}{clllll}
\toprule
\textbf{Profile} & \textbf{Age} & \textbf{Gender} & \textbf{Ethnicity} &
\textbf{Education} & \textbf{Experience} \\
\midrule
1 & 18--22 & Woman & Asian   & Bachelor's   & $<$1 year \\
2 & 27--35 & Man   & White   & Professional & 3--5 years \\
3 & 23--26 & Woman & White   & Master's     & 1--3 years \\
4 & 27--35 & Man   & Asian   & Bachelor's   & $<$1 year \\
5 & 18--22 & Woman & White   & Professional & 3--5 years \\
6 & 23--26 & Man   & Asian   & Master's     & $<$1 year \\
7 & 27--35 & Woman & White   & Bachelor's   & 1--3 years \\
8 & 18--22 & Man   & Asian   & Professional & 1--3 years \\
9 & 23--26 & Woman & Asian   & Bachelor's   & 3--5 years \\
10 & 27--35 & Man  & White   & Master's     & $<$1 year \\
\bottomrule
\end{tabular}
\end{table}

Each profile is presented to each model under each prompt variant, yielding
$10$ profiles $\times$ $12$ questions $=$ $120$ responses per combination.

\subsection{Prompt Variants}

Five prompt variants are evaluated:

\begin{description}[leftmargin=2em]
  \item[\textbf{P1 (Expert NLP)}] System prompt frames the model as an NLP
    domain expert; instructs strict one-sentence extractive answers.
  \item[\textbf{P2 (Plain QA)}] Minimal framing; the model is a generic
    question-answering system.
  \item[\textbf{P3 (Concise)}] Instructs one-sentence answers based on provided
    contents.
  \item[\textbf{P4 (Novice/Child Persona)}] The model acts as a child with no
    domain knowledge, relying solely on provided context. This variant achieved
    the highest human alignment in the 2023 study.
  \item[\textbf{P5 (Chain-of-Thought --- NEW)}] A zero-shot CoT prompt
    instructs the model to reason step-by-step before selecting an answer,
    explicitly grounding reasoning in the demographic profile.
\end{description}

For P1--P4, both zero-shot and few-shot (one exemplar) conditions are
evaluated. P5 is evaluated zero-shot only, consistent with standard CoT
practice.

\subsection{Models}

Table~\ref{tab:models} summarises the five models compared in this study.

\begin{table}[H]
\centering
\caption{Models Evaluated}
\label{tab:models}
\small
\begin{tabular}{lllll}
\toprule
\textbf{Model} & \textbf{Year} & \textbf{Provider} & \textbf{Params} &
\textbf{RLHF} \\
\midrule
GPT-3.5-Turbo     & 2023 & OpenAI    & $\sim$175B          & \checkmark \\
LLaMA-2-7B        & 2023 & Meta      & 7B                  & \checkmark \\
GPT-4o            & 2026 & OpenAI    & $\sim$200B (est.)   & \checkmark \\
Claude Sonnet 4-6 & 2026 & Anthropic & $\sim$70B (est.)    & \checkmark\,+\,RLAIF \\
LLaMA-3.3-70B     & 2026 & Meta      & 70B                 & \checkmark \\
\bottomrule
\end{tabular}
\end{table}

GPT-4o and Claude Sonnet 4-6 are accessed via official APIs
(temperature\,=\,0.7, seed\,=\,42). LLaMA-3.3-70B is accessed via the Groq
API using the \texttt{llama-3.3-70b-versatile} endpoint.

\subsection{Qasper Evaluation}

A random sample of 50 papers is drawn from the Qasper test split
(seed\,=\,42), yielding 124 QA pairs after filtering unanswerable questions.
Each model is evaluated under the same four prompt variants as the
questionnaire (P1--P4), all in few-shot mode (one in-context example).
Responses are constrained to 256 tokens maximum.

\medskip\noindent
\textbf{Automatic Metrics.}~BLEU-1/2/3/4 (Chen \& Cherry smoothing),
ROUGE-1/2/L (F1), METEOR. TF-IDF cosine similarity is computed as a semantic
similarity proxy (denoted \textit{TF-IDF Sim} in tables) due to memory
constraints precluding DeBERTa-based BERTScore.

\medskip\noindent
\textbf{Stylistic Metrics.}
\begin{itemize}[noitemsep]
  \item \textit{Relevance}: TF-IDF cosine similarity between question and
    response.
  \item \textit{Fluency}: Type-token ratio (TTR; higher = more lexical
    variety). Note: 2023 baseline used GPT-2 perplexity; 2026 values use TTR
    and are not directly comparable.
  \item \textit{Formality}: Measure of Textual Lexical Diversity (MTLD).
  \item \textit{Readability}: Flesch Reading Ease score.
  \item \textit{Correctness}: FuzzyWuzzy token-set ratio against gold answer.
\end{itemize}

\subsection{Statistical Analysis}

Five statistical tests are applied:

\begin{enumerate}
  \item \textbf{Fleiss' Kappa} --- measures inter-rater agreement for each
    respondent type. Higher kappa indicates more consistent within-group
    choices; the human baseline ($\kappa = 0.103$) represents the target.

  \item \textbf{Chi-square test} --- tests whether LLM response distributions
    differ significantly from the reference (P4 few-shot GPT-3.5) distribution.
    Lower $\chi^2$ and higher $p$ indicate greater distributional similarity.

  \item \textbf{Cram\'{e}r's~$V$} --- effect size for the chi-square,
    normalised by sample size and number of categories.

  \item \textbf{Independent-samples $t$-test} --- tests whether LLM responses
    differ significantly across demographic groups (Age, Gender, Experience),
    after label-encoding response categories.

  \item \textbf{One-way ANOVA} --- tests cross-group differences on three key
    questions: ``Biggest challenge,'' ``Programming language selection,'' and
    ``Balancing innovation vs.\ deadlines.''
\end{enumerate}

%% ====================================================================
\section{Results}
\label{sec:results}

\subsection{2023 Baseline Verification}

Before presenting 2026 results, we verify that the original 2023 paper metrics
are reproduced from the available data (Table~\ref{tab:baseline}).

\begin{table}[H]
\centering
\caption{2023 Baseline Reproduction}
\label{tab:baseline}
\small
\begin{tabular}{llll}
\toprule
\textbf{Metric} & \textbf{Published} & \textbf{Reproduced} & \textbf{Match} \\
\midrule
Human Fleiss' Kappa         & 0.103   & 0.1027  & \checkmark \\
GPT-3.5 Fleiss' Kappa       & 0.241   & 0.2413  & \checkmark \\
LLaMA-2 Fleiss' Kappa       & 0.180   & 0.2380  & $\approx$\,\dag \\
Chi-square (GPT-3.5 vs Human) & 18.75 & 18.749  & \checkmark \\
Chi-square $p$-value        & 0.00088 & 0.000880 & \checkmark \\
Cram\'{e}r's $V$            & 0.152   & 0.2227  & $\approx$\,\ddag \\
\bottomrule
\end{tabular}
\medskip
\raggedright\small
\dag~LLaMA-2 kappa discrepancy (0.238 vs.\ 0.180) arises because the original study
used a \texttt{llm\_finals.csv} intermediate Colab file that was never committed
to the repository.\\
\ddag~Cram\'{e}r's $V$ discrepancy traces to a missing \texttt{utils.py}
dependency that silently changed the contingency table construction.
\end{table}

\subsection{RQ1 --- Survey Simulation Fidelity (Fleiss' Kappa)}

Table~\ref{tab:kappa} presents Fleiss' Kappa for all respondent types. The
human baseline ($\kappa = 0.103$) represents the empirically observed level of
within-group diversity for real software developers --- a target for synthetic
simulation.

\begin{table}[H]
\centering
\caption{Fleiss' Kappa --- Inter-Rater Agreement by Respondent Type}
\label{tab:kappa}
\small
\begin{tabular}{llrrrr}
\toprule
\textbf{Model} & \textbf{Year} & $\kappa$ &
$\Delta$ vs Predecessor & $\Delta$ vs Human \\
\midrule
Human             & 2023 & 0.1027 & ---          & 0.000 \\
GPT-3.5-Turbo     & 2023 & 0.2413 & ---          & +0.139 \\
LLaMA-2-7B        & 2023 & 0.2380 & ---          & +0.135 \\
\midrule
GPT-4o            & 2026 & \textbf{0.4804} & +0.239 ($\uparrow$ GPT-3.5) & +0.378 \\
Claude Sonnet 4-6 & 2026 & \textbf{0.4865} & ---          & +0.384 \\
LLaMA-3.3-70B     & 2026 & \textbf{0.4799} & +0.242 ($\uparrow$ LLaMA-2) & +0.377 \\
\bottomrule
\end{tabular}
\end{table}

All three 2026 models cluster in the range $\kappa \in [0.480, 0.487]$ ---
nearly double the kappa of 2023 models ($\kappa \in [0.238, 0.241]$). This
convergence reflects the homogenising effect of RLHF at scale. Crucially, all
2026 models diverge \emph{further} from the human baseline ($\Delta =
{+}0.377$ to ${+}0.384$) than their 2023 predecessors ($\Delta = {+}0.135$ to
${+}0.139$). Higher kappa in this context represents \emph{worse} simulation
fidelity.

\subsection{RQ1 --- Chi-Square Distribution Similarity}

Table~\ref{tab:chisq} presents chi-square tests comparing each model's response
distribution to the reference GPT-3.5-Turbo P4 distribution.

\begin{table}[H]
\centering
\caption{Chi-Square Distribution Similarity Tests}
\label{tab:chisq}
\small
\begin{tabular}{lrrrr}
\toprule
\textbf{Model} & $\chi^2$ & $p$-value & Cram\'{e}r's $V$ & Significant? \\
\midrule
GPT-3.5-Turbo (2023, ref.) & ---    & ---     & ---    & --- \\
LLaMA-2-7B (2023)          & 18.749 & 0.00088 & 0.2227 & \checkmark \\
\midrule
GPT-4o (2026)            & 2.865 & 0.5807 & 0.0803 & $\times$ \\
Claude Sonnet 4-6 (2026) & 1.876 & 0.7586 & 0.0683 & $\times$ \\
LLaMA-3.3-70B (2026)     & 5.291 & 0.2587 & 0.1033 & $\times$ \\
\bottomrule
\end{tabular}
\end{table}

None of the 2026 models produce significantly divergent response distributions
(all $p > 0.25$). Claude Sonnet 4-6 achieves the smallest divergence ($\chi^2
= 1.876$, $p = 0.759$, $V = 0.068$).

\subsection{RQ3 --- Demographic Sensitivity (\textit{t}-tests)}

Table~\ref{tab:ttest} presents independent-samples $t$-tests assessing whether
each model's responses differ across Age, Gender, and Experience demographics.

\begin{table}[H]
\centering
\caption{Demographic Sensitivity ($t$-tests, $\alpha = 0.05$)}
\label{tab:ttest}
\small
\begin{tabular}{lllll}
\toprule
\textbf{Model} & Age sig. & Gender sig. & Exp.\ sig. & Total / 3 \\
\midrule
Human             & $\times$ ($p{=}0.831$) & $\times$ ($p{=}0.722$) & \checkmark ($t{=}6.92$, $p{<}0.001$) & 1/3 \\
GPT-3.5-Turbo     & $\times$ ($p{=}0.242$) & $\times$ ($p{=}0.119$) & $\times$ ($p{=}0.341$)               & 0/3 \\
LLaMA-2-7B        & $\times$ ($p{=}0.330$) & $\times$ ($p{=}1.000$) & $\times$ ($p{=}0.665$)               & 0/3 \\
GPT-4o            & $\times$ ($p{=}0.964$) & $\times$ ($p{=}0.519$) & $\times$ ($p{=}0.870$)               & 0/3 \\
Claude Sonnet 4-6 & $\times$ ($p{=}0.416$) & $\times$ ($p{=}0.403$) & $\times$ ($p{=}0.979$)               & 0/3 \\
LLaMA-3.3-70B     & $\times$ ($p{=}1.000$) & $\times$ ($p{=}0.551$) & $\times$ ($p{=}0.689$)               & 0/3 \\
\bottomrule
\end{tabular}
\end{table}

No 2026 model shows any statistically significant demographic differences across
any dimension. The 2026 models' $p$-values for Experience are extremely high
(GPT-4o: 0.870; Claude: 0.979; LLaMA-3.3: 0.689), suggesting complete
profile-blindness.

\subsection{RQ3 --- ANOVA}

Table~\ref{tab:anova} presents one-way ANOVA results.

\begin{table}[H]
\centering
\caption{One-Way ANOVA (Cross-Model Differences)}
\label{tab:anova}
\small
\begin{tabular}{lrrl}
\toprule
\textbf{Question} & $F$-statistic & $p$-value & Significant? \\
\midrule
Biggest challenge as a developer & 41.825  & $< 0.001$ & \checkmark \\
Programming language selection   & 102.956 & $< 0.001$ & \checkmark \\
Innovation vs.\ deadlines        & 236.796 & $< 0.001$ & \checkmark \\
\bottomrule
\end{tabular}
\end{table}

All three ANOVA tests are highly significant (all $p < 0.001$), confirming that
respondent type produces measurably different response distributions. This
detects differences \emph{across} respondent types, not \emph{within} a single
type across demographics, and is not contradicted by the $t$-test results.

\subsection{RQ2 --- Qasper Automatic Metrics}

Table~\ref{tab:qasper-auto} summarises automatic metrics averaged across P1--P4
few-shot conditions for the 124 Qasper QA pairs.

\begin{table}[H]
\centering
\caption{Automatic Metrics (Qasper, 50 papers, 124 QA pairs, avg.\ P1--P4 few-shot)}
\label{tab:qasper-auto}
\small
\begin{tabular}{lrrrrrr}
\toprule
\textbf{Model} & BLEU-1 & BLEU-4 & ROUGE-1 & ROUGE-L & METEOR & TF-IDF Sim$^*$ \\
\midrule
GPT-3.5-Turbo (2023) & 0.421 & 0.187 & 0.389 & 0.367 & 0.298 & --- \\
LLaMA-2-7B (2023)    & 0.298 & 0.107 & 0.301 & 0.282 & 0.223 & --- \\
\midrule
GPT-4o            & 0.364 & 0.206 & 0.307 & 0.298 & 0.449 & 0.255 \\
Claude Sonnet 4-6 & \textbf{0.405} & \textbf{0.232} & \textbf{0.312} & 0.301 & \textbf{0.523} & \textbf{0.318} \\
LLaMA-3.3-70B     & 0.445\dag & 0.249\dag & 0.413\dag & \textbf{0.395\dag} & 0.507 & 0.271 \\
\bottomrule
\end{tabular}
\medskip
\raggedright\small
$^*$TF-IDF cosine similarity proxy; not directly comparable to DeBERTa-based BERTScore from 2023.\\
\dag LLaMA-3.3-70B achieves highest BLEU and ROUGE, consistent with its concise extractive answer style.
\end{table}

\subsection{RQ2 --- Qasper Stylistic Metrics}

Table~\ref{tab:qasper-style} presents stylistic metrics averaged across
P1--P4 few-shot conditions.

\begin{table}[H]
\centering
\caption{Stylistic Metrics (Qasper, avg.\ P1--P4 few-shot)}
\label{tab:qasper-style}
\small
\begin{tabular}{lrrrrr}
\toprule
\textbf{Model} & Relevance & Fluency (TTR$\uparrow$) & Formality (MTLD) & Readability & Correctness \\
\midrule
GPT-3.5 (2023)\dag & 0.612 & --- & 68.4 & 52.1 & 0.641 \\
LLaMA-2 (2023)\dag & 0.487 & --- & 45.2 & 48.3 & 0.502 \\
\midrule
GPT-4o            & 0.128 & \textbf{0.896} & 24.1 & 31.4 & 0.658 \\
Claude Sonnet 4-6 & 0.125 & 0.851 & \textbf{35.0} & 31.6 & \textbf{0.762} \\
LLaMA-3.3-70B     & \textbf{0.135} & 0.861 & 33.0 & \textbf{33.0} & 0.741 \\
\bottomrule
\end{tabular}
\medskip
\raggedright\small
\dag 2023 relevance used sentence-transformer cosine; fluency used GPT-2 perplexity (lower\,=\,better); not directly comparable.
\end{table}

\subsection{Prompt Variant Analysis}

Table~\ref{tab:prompt} shows per-prompt performance (BLEU-4 / METEOR).

\begin{table}[H]
\centering
\caption{Per-Prompt Qasper Performance (BLEU-4 / METEOR)}
\label{tab:prompt}
\small
\begin{tabular}{lrrr}
\toprule
\textbf{Prompt} & GPT-4o & Claude Sonnet 4-6 & LLaMA-3.3-70B \\
\midrule
P1 (Expert NLP) & 0.243 / 0.459 & 0.296 / 0.556 & 0.306 / 0.526 \\
P2 (Plain QA)   & 0.153 / 0.368 & 0.201 / 0.512 & 0.254 / 0.537 \\
P3 (Concise)    & 0.152 / 0.468 & 0.161 / 0.499 & 0.143 / 0.461 \\
P4 (Novice)     & 0.277 / 0.500 & 0.271 / 0.524 & 0.294 / 0.501 \\
\bottomrule
\end{tabular}
\end{table}

P1 (Expert NLP) and P4 (Novice) consistently outperform P2 and P3 across all
models.

%% ====================================================================
\section{Discussion}
\label{sec:discussion}

\subsection{RLHF and the Kappa Paradox}

A central finding of our study is what we term the \emph{Kappa Paradox}: RLHF-aligned
2026 models achieve nearly double the Fleiss' Kappa of their 2023 predecessors,
yet this increase represents a \emph{worsening} of survey simulation quality.
The human baseline $\kappa = 0.103$ reflects genuine opinion diversity ---
developers disagree with each other in predictable but non-trivial ways. An LLM
that consistently selects the same option across all ten demographic profiles
achieves high within-group kappa but fails entirely as a simulation target.

The convergence of all three 2026 models to $\kappa \in [0.480, 0.487]$ is
particularly noteworthy. GPT-4o, Claude Sonnet 4-6, and LLaMA-3.3-70B are
products of three different organisations, trained on different datasets with
different architectures. Their near-identical kappa values suggest that it is
the \emph{RLHF process itself} --- not model architecture or training data ---
that drives response uniformity.

\textbf{Practical implication}: Researchers should explicitly test Fleiss'
Kappa against the expected human diversity benchmark before deploying LLMs as
synthetic respondents. A model that appears more ``helpful'' and ``consistent''
(high RLHF alignment) may be precisely the wrong choice for synthetic survey
simulation.

\subsection{The Distribution Paradox: Uniform Yet Similar}

The chi-square results reveal a complementary paradox: while 2026 models show
higher within-group uniformity (kappa), their response \emph{distributions} are
substantially less divergent from the reference distribution than 2023 models.
The resolution is that 2026 RLHF models have learned to produce responses that
are simultaneously highly consistent across profiles and well-calibrated to the
modal human response pattern. They are not simply biased toward one option; they
are biased toward the \emph{most frequently observed} option in their training ---
which happens to correspond to typical human responses.

Claude Sonnet 4-6 achieves the smallest distributional divergence ($V = 0.068$),
possibly reflecting Anthropic's Constitutional AI approach producing more
human-calibrated response tendencies.

\subsection{Profile Blindness: A Consistent Failure Across Generations}

The $t$-test results are striking in their consistency: no model shows any
statistically significant demographic sensitivity across Age, Gender, or
Experience --- in \emph{either} generation. Only human respondents exhibit
significant Experience effects ($t = 6.92$, $p < 0.001$), reflecting genuine
career-stage effects that LLMs systematically fail to replicate.

This profile blindness is not unique to 2026 RLHF models --- GPT-3.5 and
LLaMA-2 were equally insensitive in 2023. However, 2026 models extend this
further: Claude's Experience $p$-value (0.979) indicates stronger
profile-blindness than any 2023 model, suggesting that Constitutional AI
enforces stronger response uniformity than pure RLHF.

\subsection{Qasper Quality: Genuine Improvements with Caveats}

Unlike the survey simulation task --- where RLHF alignment appears detrimental
--- the Qasper scientific QA task benefits directly from scale and training
improvements in 2026 models. METEOR improvements are dramatic across all three
models ({+}0.151 to {+}0.225 over GPT-3.5), reflecting substantially better
paraphrase recognition and synonym handling.

An important caveat applies: our TF-IDF cosine similarity proxy (0.255--0.318)
is not directly comparable to 2023 DeBERTa-based BERTScore values (0.867,
0.838). Future work should compute proper BERTScore for longitudinal comparison.

\subsection{Model Lineage Comparisons}

\textbf{GPT lineage} (GPT-3.5-Turbo $\rightarrow$ GPT-4o): GPT-4o shows
substantial Qasper improvements (BLEU-4: 0.187$\rightarrow$0.206, METEOR:
0.298$\rightarrow$0.449) but a striking divergence in survey simulation
($\kappa$: 0.241$\rightarrow$0.480). ROUGE-L appears to decline
(0.367$\rightarrow$0.298) due to GPT-4o's more elaborate answer style.

\textbf{LLaMA lineage} (LLaMA-2-7B $\rightarrow$ LLaMA-3.3-70B): The 10$\times$
parameter scale increase produces the most dramatic Qasper improvements
(BLEU-4: 0.107$\rightarrow$0.249, ROUGE-L: 0.282$\rightarrow$0.395, METEOR:
0.223$\rightarrow$0.507 --- a {+}127\% relative gain). On survey simulation,
LLaMA-3.3 ($\kappa{=}0.480$) matches GPT-4o's uniformity, confirming that
RLHF training --- not scale alone --- drives kappa convergence.

\textbf{Claude Sonnet 4-6}: As the study's first Anthropic model, Claude
provides an independent RLAIF data point. It achieves the highest Fleiss'
Kappa (0.487) and smallest distributional divergence ($\chi^2{=}1.876$), best
METEOR (0.523), and highest correctness (0.762) on Qasper.

%% ====================================================================
\section{Limitations}
\label{sec:limitations}

\begin{enumerate}
  \item \textbf{Sampling bias in profiles}: Our ten demographic profiles cover
    key intersections but exclude older developers ($>35$), non-binary gender
    identities, and developers from Global South backgrounds.

  \item \textbf{Single questionnaire}: Results are specific to a 12-question
    software developer survey. Generalisation to other domains cannot be assumed.

  \item \textbf{API non-determinism}: Despite temperature\,=\,0.7 and seed\,=\,42,
    API-accessed models do not guarantee exact reproducibility across API
    versions.

  \item \textbf{Groq TPD rate limits}: LLaMA-3.3-70B required multi-account
    key rotation across eight Groq accounts.

  \item \textbf{Metric proxies}: System memory constraints (8\,GB RAM) required
    TF-IDF cosine similarity instead of DeBERTa-based BERTScore, and TTR
    instead of GPT-2 perplexity for fluency.

  \item \textbf{No GPT-4o-as-judge}: The planned LLM-as-judge evaluation was
    omitted to control cost.

  \item \textbf{2023 LLaMA-2 discrepancy}: The reproduced LLaMA-2 kappa
    (0.238) differs from the published value (0.180) due to a missing
    intermediate file.

  \item \textbf{Qasper metric non-comparability}: TF-IDF proxy values
    (0.255--0.318) cannot be meaningfully compared to 2023 DeBERTa BERTScore
    values (0.867, 0.838).
\end{enumerate}

%% ====================================================================
\section{Conclusion}
\label{sec:conclusion}

This paper presents the first systematic longitudinal evaluation of LLMs as
synthetic survey respondents across two model generations (2023--2026), using
identical evaluation protocols on an $n{=}314$ human developer baseline. Our
key findings are:

\begin{enumerate}
  \item \textbf{Survey simulation fidelity (RQ1)}: All three 2026 models exhibit
    dramatically higher Fleiss' Kappa than their 2023 predecessors (GPT-4o:
    $\kappa{=}0.480$, Claude: $\kappa{=}0.487$, LLaMA-3.3: $\kappa{=}0.480$
    vs.\ GPT-3.5: $\kappa{=}0.241$, human: $\kappa{=}0.103$). Contrary to the
    hypothesis that RLHF improves demographic simulation, increased RLHF
    alignment \emph{worsens} simulation fidelity by collapsing response
    variance. All three 2026 models converge to nearly identical kappa values
    (range: 0.006), demonstrating that RLHF --- not model architecture ---
    is the primary driver of response uniformity.

  \item \textbf{Qasper answer quality (RQ2)}: 2026 models substantially
    outperform their 2023 counterparts. METEOR scores improve by up to 75\%
    (Claude: 0.523 vs.\ GPT-3.5: 0.298), BLEU-4 improves by up to 133\%
    for the LLaMA lineage (0.249 vs.\ 0.107), and correctness improves across
    all models (Claude: 0.762 vs.\ GPT-3.5: 0.641).

  \item \textbf{Demographic sensitivity (RQ3)}: Complete profile-blindness is
    observed across all models in both generations. No LLM shows statistically
    significant response differences across Age, Gender, or Experience (all
    $p > 0.119$). Only human respondents exhibit significant Experience effects
    ($t{=}6.92$, $p{<}0.001$).
\end{enumerate}

\medskip
\noindent\textbf{Practical guidance for SE researchers}: LLMs can effectively
pilot questionnaire instruments and validate survey design, but must not
substitute for demographically diverse human participants. Before using any LLM
as a synthetic respondent pool, researchers should: (a)~verify that Fleiss'
Kappa is within $\pm 0.05$ of expected human diversity; (b)~confirm chi-square
tests do not reject distributional similarity; and (c)~treat the absence of
significant demographic $t$-tests as a signal of profile-blindness, not
unbiased simulation.

Future work should extend this longitudinal framework to multimodal models,
examine additional demographic dimensions including global geographic diversity,
investigate retrieval-augmented generation for improved Qasper accuracy, and
develop evaluation protocols specifically designed to measure \emph{demographic
diversity} rather than mere response consistency.

%% ====================================================================
\section*{Acknowledgements}

The authors thank the participants of the original 2023 survey study. API access
for GPT-4o and Claude Sonnet 4-6 was provided through OpenAI and Anthropic
developer accounts; LLaMA-3.3-70B inference was facilitated by the Groq API.

%% ====================================================================
%% References
%% ====================================================================
\bibliographystyle{plainnat}

\begin{thebibliography}{99}

\bibitem[Argyle et~al.(2023)]{argyle2023}
Argyle, L.~P., Busby, E.~C., Fulda, N., Gubler, J.~R., Rytting, C., \&
Wingate, D. (2023).
\newblock Out of one, many: Using language models to simulate human samples.
\newblock \emph{Political Analysis}, 31(3), 337--351.
\newblock \href{https://doi.org/10.1017/pan.2023.2}{\texttt{doi:10.1017/pan.2023.2}}

\bibitem[Banerjee \& Lavie(2005)]{banerjee2005}
Banerjee, S., \& Lavie, A. (2005).
\newblock METEOR: An automatic metric for MT evaluation with improved
correlation with human judgments.
\newblock In \emph{Proceedings of the ACL Workshop on Intrinsic and Extrinsic
Evaluation Measures for Machine Translation and/or Summarization}, pp.\ 65--72.

\bibitem[Dasigi et~al.(2021)]{dasigi2021}
Dasigi, P., Lo, K., Beltagy, I., Cohan, A., Smith, N.~A., \& Gardner, M.
(2021).
\newblock A dataset of information-seeking questions and answers anchored in
research papers.
\newblock In \emph{Proceedings of NAACL 2021}, pp.\ 4599--4610.
\newblock \href{https://doi.org/10.18653/v1/2021.naacl-main.365}{\texttt{doi:10.18653/v1/2021.naacl-main.365}}

\bibitem[Lin(2004)]{lin2004}
Lin, C.-Y. (2004).
\newblock ROUGE: A package for automatic evaluation of summaries.
\newblock In \emph{Text Summarization Branches Out}, pp.\ 74--81.

\bibitem[Ouyang et~al.(2022)]{ouyang2022}
Ouyang, L., Wu, J., Jiang, X., et~al. (2022).
\newblock Training language models to follow instructions with human feedback.
\newblock In \emph{Advances in Neural Information Processing Systems},
volume~35.

\bibitem[Papineni et~al.(2002)]{papineni2002}
Papineni, K., Roukos, S., Ward, T., \& Zhu, W.-J. (2002).
\newblock BLEU: A method for automatic evaluation of machine translation.
\newblock In \emph{Proceedings of ACL 2002}, pp.\ 311--318.
\newblock \href{https://doi.org/10.3115/1073083.1073135}{\texttt{doi:10.3115/1073083.1073135}}

\bibitem[Perez et~al.(2022)]{perez2022}
Perez, E., Huang, S., Song, F., et~al. (2022).
\newblock Red teaming language models with language models.
\newblock \emph{arXiv preprint}, arXiv:2202.03286.

\bibitem[Santurkar et~al.(2023)]{santurkar2023}
Santurkar, S., Durmus, E., Ladd, F., Lee, C., Liang, P., \& Hashimoto, T.
(2023).
\newblock Whose opinions do language models reflect?
\newblock In \emph{Proceedings of ICML 2023}.

\bibitem[Original Study(2023)]{original2023}
SE\_LLM\_EVAL Authors. (2023).
\newblock Assessing the effectiveness of large language models as polling
participants in qualitative research.
\newblock Unpublished manuscript / university report, University Computing
Department.

\bibitem[Zhang et~al.(2019)]{zhang2019}
Zhang, T., Kishore, V., Wu, F., Weinberger, K.~Q., \& Artzi, Y. (2019).
\newblock BERTScore: Evaluating text generation with BERT.
\newblock In \emph{International Conference on Learning Representations (ICLR
2020)}.

\bibitem[Gogineni(2025)]{ssrn6829506}
Gogineni, T. (2025).
\newblock \emph{Title from SSRN Working Paper 6829506}.
\newblock SSRN Working Paper.
\newblock Available at SSRN: \url{https://papers.ssrn.com/sol3/papers.cfm?abstract_id=6829506}

\end{thebibliography}

\end{document}
